% \documentclass[a4paper,12pt]{article}
\documentclass[a4paper,11pt]{article}
% -------------------------------------------------
% Set up the page margins.
% -------------------------------------------------
% \usepackage[text={6.5in,9.5in},centering]{geometry}
% \setlength{\topmargin}{-0.25in}
\usepackage[margin=1in]{geometry}

% -------------------------------------------------
% Load Packages here
% -------------------------------------------------
%\usepackage{hyperref} 
\usepackage{graphicx,url,subfig}
\RequirePackage[OT1]{fontenc}
\RequirePackage{amsthm,amsmath,amssymb,amscd}
\RequirePackage[numbers]{natbib}
\RequirePackage[colorlinks,citecolor=blue,urlcolor=blue]{hyperref}
\usepackage{graphics}
\usepackage{enumerate}
\usepackage{datetime}
\usepackage{etex}
% \usepackage[notcite,notref,color]{showkeys}
\usepackage[normalem]{ulem} % either use this (simple) or
\usepackage{soul} % use this (many fancier options)
\makeatother
\numberwithin{equation}{section}
\allowdisplaybreaks
\usepackage{mathrsfs}


% -------------------------------------------------
% check references
% -------------------------------------------------
% \usepackage{refcheck}
% \usepackage[notref,notcite]{showkeys}
% \usepackage[notcite]{showkeys}
% \usepackage{showkeys}

% -------------------------------------------------
% Environments here
% -------------------------------------------------
\theoremstyle{plain}
\newtheorem{theorem}{Theorem}[section]
\newtheorem{corollary}[theorem]{Corollary}
\newtheorem{lemma}[theorem]{Lemma}
\newtheorem{proposition}[theorem]{Proposition}

\theoremstyle{definition}
\newtheorem{definition}[theorem]{Definition}
\newtheorem{hypothesis}[theorem]{Hypothesis}
\newtheorem{assumption}[theorem]{Assumption}

\newtheorem{remark}[theorem]{Remark}
\newtheorem{conjecture}[theorem]{Conjecture}
\newtheorem{example}[theorem]{Example}

% -------------------------------------------------
%  Define short hand symbols.
% -------------------------------------------------
\newcommand{\E}{\mathbb{E}}
\newcommand{\W}{\dot{W}}
\newcommand{\B}{\textbf{B}}
\newcommand{\D}{\mathbb{D}}
\newcommand{\ud}{\ensuremath{\mathrm{d}}}
\newcommand{\Ceil}[1]{\left\lceil #1 \right\rceil}
\newcommand{\Floor}[1]{\left\lfloor #1 \right\rfloor}
\newcommand{\sgn}{\text{sgn}}
\newcommand{\Lad}{\text{L}_{\text{ad}}^2}
\newcommand{\SI}[1]{\mathcal{I}\left[#1 \right]}
\newcommand{\SIB}[2]{\mathcal{I}_{#2}\left[#1 \right]}
\newcommand{\Indt}[1]{1_{\left\{#1 \right\}}}
\newcommand{\LadInPrd}[1]{\left\langle #1 \right\rangle_{\text{L}_\text{ad}^2}}
\newcommand{\LadNorm}[1]{\left|\left|  #1 \right|\right|_{\text{L}_\text{ad}^2}}
\newcommand{\Norm}[1]{\left\|  #1   \right\|}
\newcommand{\Ito}{It\^{o} }
\newcommand{\Itos}{It\^{o}'s }
\newcommand{\spt}[1]{\text{supp}\left(#1\right)}
\newcommand{\InPrd}[1]{\left\langle #1 \right\rangle}
\newcommand{\mr}{\textbf{r}}
\newcommand{\Ei}{\text{Ei}}
\newcommand{\arctanh}{\operatorname{arctanh}}
\newcommand{\ind}[1]{\mathbb{I}_{\left\{ {#1} \right\} }}

\DeclareMathOperator{\esssup}{\ensuremath{ess\,sup}}

\newcommand{\steps}[1]{\vskip 0.3cm \textbf{#1}}
\newcommand{\calB}{\mathcal{B}}
\newcommand{\calD}{\mathcal{D}}
\newcommand{\calE}{\mathcal{E}}
\newcommand{\calF}{\mathcal{F}}
\newcommand{\calG}{\mathcal{G}}
\newcommand{\calK}{\mathcal{K}}
\newcommand{\calH}{\mathcal{H}}
\newcommand{\calI}{\mathcal{I}}
\newcommand{\calL}{\mathcal{L}}
\newcommand{\calM}{\mathcal{M}}
\newcommand{\calN}{\mathcal{N}}
\newcommand{\calO}{\mathcal{O}}
\newcommand{\calT}{\mathcal{T}}
\newcommand{\G}{\mathcal{G}}
\newcommand{\calP}{\mathcal{P}}
\newcommand{\calR}{\mathcal{R}}
\newcommand{\calS}{\mathcal{S}}
\newcommand{\bbC}{\mathbb{C}}
\newcommand{\bbN}{\mathbb{N}}
\newcommand{\bbP}{\mathbb{P}}
\newcommand{\bbZ}{\mathbb{Z}}
\newcommand{\myVec}[1]{\overrightarrow{#1}}
\newcommand{\sincos}{\begin{array}{c} \cos \\ \sin \end{array}\!\!}
\newcommand{\CvBc}[1]{\left\{\:#1\:\right\}}
\newcommand*{\one}{{{\rm 1\mkern-1.5mu}\!{\rm I}}}
\newcommand{\RR}{\mathbb{R}}
\newcommand{\R}{\mathbb{R}}
\newcommand{\myEnd}{\hfill$\square$}
\newcommand{\ds}{\displaystyle}
\newcommand{\Shi}{\text{Shi}}
\newcommand{\Chi}{\text{Chi}}
\newcommand{\Daw}{\ensuremath{\mathrm{daw}}}
\newcommand{\Erf}{\ensuremath{\mathrm{erf}}}
\newcommand{\Erfi}{\ensuremath{\mathrm{erfi}}}
\newcommand{\Erfc}{\ensuremath{\mathrm{erfc}}}
\newcommand{\He}{\ensuremath{\mathrm{He}}}

\def\crtL{\theta}

\newcommand{\conRef}[1]{(#1)}

\DeclareMathOperator{\Lip}{\mathit{L}}
\DeclareMathOperator{\LIP}{Lip}
\DeclareMathOperator{\lip}{\mathit{l}}

\DeclareMathOperator{\Vip}{\overline{\varsigma}}
\DeclareMathOperator{\vip}{\underline{\varsigma}}
\DeclareMathOperator{\vv}{\varsigma}

% % Set up mdframe for questions.
\newcounter{NQ}
\newcounter{LQ}
\newcounter{ZQ}
\usepackage{xcolor}
\usepackage[framemethod=tikz]{mdframed}
\mdfsetup{%
	middlelinecolor=lightgray,
	middlelinewidth=2pt,
	backgroundcolor=red!20,
	roundcorner=10pt}
\newenvironment{NickQuestion}[1][] %
{\refstepcounter{NQ}
	\begin{mdframed}[backgroundcolor=lightgray]\textbf{Nick's Question \theNQ \: #1~~}~~~\noindent}%
	{\end{mdframed}}
\newenvironment{LeQuestion}[1][] %
{\refstepcounter{LQ}
	\begin{mdframed}[backgroundcolor=red!10]\textbf{Le's Question \theLQ \: #1~~}~~~\noindent}%
	{\end{mdframed}}
\newenvironment{ZhijianQuestion}[1][] %
{\refstepcounter{ZQ}
	\begin{mdframed}[backgroundcolor=blue!10]\textbf{Zhijian's Question \theZQ \: #1~~}~~~\noindent}%
	{\end{mdframed}}
\numberwithin{NQ}{section}
\numberwithin{LQ}{section}
\numberwithin{ZQ}{section}

\title{Nick Eisenberg's Meeting Notes}
\author{Le Chen, Nick Eisenberg, Zhijian Wu}
\date{October 2019}

\usepackage{natbib}
\usepackage{graphicx}


\begin{document}
	
\section{Chapter 1}

\begin{NickQuestion}
	Proof of lemma 1.8.4

	\textcolor{magenta}{The proof is very nice.} 
\end{NickQuestion}

\section{Chapter 2}
\begin{NickQuestion}
	In the proof of Theorem 2.4.3 on the bottom of page 34 there are some claims made that I am unable to verify why they are true. First, 
	we are considering 
	\[
	X_t = X_0 + \int_0^t u(s)dB_s
	\]
	and here $u(s)$ only satisfies 
	\[
	\mathbb P \left(\int_0^t u(s)^s ds < \infty \right) = 1
	\]
	so, as in the Kuo text, this stochastic integral exists and a limit in probability of a sequence of stochastic integrals $\int_0^t u_n(s)dB_s$ where the $u_n$ are simple and are in $L^2_{ad}$ and hence $\int_0^t u_n(s)dB_s$ is a martingale and we know the first and second moment. But for the limit $ \int_0^t u(s)dB_s$ we know nothing about the moments. But in this text, it is claimed that 
	\[
	\E \left[\int_{t_i}^{t_{i+1}}u^2(s)ds - \left(\int_{t_i}^{t_{i+1}}u(s)dB_s\right)^2 \bigg| \mathcal F_i\right] =0
	\]
	and I have never seen this claimed in any other text. 

	\textcolor{magenta}{Note that $u$ is bounded (we are assuming localization has been done). Hence, all moments exists.} 
\end{NickQuestion}


\begin{NickQuestion}
	In the proof of theorem 2.5.1, why is it true that 
	\[ 
		E \left(\int_0^t |u_n'(B_s)-1_{(x,\infty)}(B_s)|^2 ds\right) \le \int_0^t \mathbb P \left( |B_s-x|\le \frac{2}{n}\right)ds
	\]
	\textcolor{magenta}{Please explore the mollifier $f$ and $f_n$. Plot these $u_n'$ and see how could it differ from $1_{(x,\infty)}$.} 
\end{NickQuestion}	

\begin{NickQuestion}
	What is the motivation on page 41 behind the following definition of the backward ito integral 
	\[
	\int_0^T u(s)\hat{d}B_s = \lim_{n \to \infty} \sum_{i=0}^{n-2}\left( \frac{n}{T}\int_{(j+1)T/n}^{(j+2)T/n}u(s)ds \right)(B_{(j+1)T/n}-B_{jT/n})
	\]
	The integral in the parenthesis is the average of $u_s$ on $\big((j+1)T/n,(j+2)T/n\big)$. 

	\bigskip
	\textcolor{magenta}{Let's skip this part. I know nothing on the backward SDE and not particular interested either. 
	It has applications in finance. If you don't do finance, I don't see the point here.} 
\end{NickQuestion}

\begin{NickQuestion}[(Problem 2.7)]
	This problem follows easily from Levy's Characterization of Brownian Motion. But I can't figure out how to do this problem with the hint given. 

	\bigskip
	\textcolor{magenta}{This is a nice problem. The hints is to ask you to find the distribution of $X_t-X_s$ 
	through the characteristic function. Once apply Ito's formula, take expectation and remove the martingale part, you should 
	arrive at an integral equation. Take derivatives, you arrive at a differential equation. The solution for $a_s\left(t\right):=\E\left(e^{i\lambda \left(X_t-X_s\right)}\right)$
	is 
	\begin{align*}
		a_s(t) = e^{-\frac{\lambda^2}{2}(t-s)}.
	\end{align*}
	} 
\end{NickQuestion}

\begin{NickQuestion}[(Problem 2.15)]
	First we consider the $2-d$ Ito process
	\[
		\begin{pmatrix}
			M_s 
			\\ N_s
		\end{pmatrix} 
		=
		\begin{pmatrix}
		E(X^p) 
		\\ E(Y^{q'})
		\end{pmatrix}
		+
	\int_0^t {	\begin{pmatrix}
		\varphi_s & 0
		\\0 & \psi_s
		\end{pmatrix}
		\begin{pmatrix}
		\hat{d B_s} 
		\\ d B_s
		\end{pmatrix} }
	\]
	and then following the hint we apply Ito's formula to $f(x,y) = x^{1/p}y^{1/q'}$. We get the following derivatives:
	\begin{itemize}
		\item $f_x = \frac{1}{p}x^{\frac{1}{p}-1}y^{\frac{1}{q'}}$
		\item $f_{xx} = \frac{1}{p} \left(\frac{1}{p}-1\right)x^{\frac{1}{p}-2}y^{\frac{1}{q'}}$
		\item $f_y = \frac{1}{q'}x^{\frac{1}{p}}y^{\frac{1}{q'}-1}$
		\item $f_{yy} =\frac{1}{q'} \left(\frac{1}{q'}-1\right)x^{\frac{1}{p}}y^{\frac{1}{q'}-2}$
	\end{itemize} 
	Now putting this into Ito's Formula gives us 
	\begin{align*}
		M_t^{1/p}N_t^{1/q'}&=\Norm{X}_p\Norm{Y}_{q'} + \frac{1}{p}\int_0^t M_s^{\frac{1}{p}-1}N_s^{\frac{1}{q'}} \varphi_s \: d\hat{B_s}+ \frac{1}{q'}\int_0^t M_s^{\frac{1}{p}}N_s^{\frac{1}{q'}-1}\psi_s \: dB_s
		\\& + \frac{1}{2}\frac{1}{p} \left(\frac{1}{p}-1\right)\int_0^t M_s^{\frac{1}{p}-2}N_s^{\frac{1}{q'}}\varphi_s^2 \: ds + \frac{1}{2}\frac{1}{q'} \left(\frac{1}{q'}-1\right)\int_0^t M_s^{\frac{1}{p}}N_s^{\frac{1}{q'}-2} \psi_s^2 \: ds
	\end{align*}

	I am not really sure where to continue from here. I see that if we let $t \to \infty$ then $M_t^{1/p}N_t^{1/q'} \to XY$ so if we took the expected value of both sides then we would partially have what we are looking for since $E(XY)$ is on the left and $\Norm{X}_p\Norm{Y}_{q'}$ is on the right. But I am not sure about the other terms. Im guessing the Ito integrals in the right have $0$ expectation but I am not sure if that is in fact true since I dont see how it is easily seen that the integrands are in $L^2(\mathcal P)$. 

	What I did next was see what $\hat{B_s}$ was equal too and for this we define another Ito Process
		\[
	\begin{pmatrix}
	B_s 
	\\ \tilde{B_s}
	\end{pmatrix} 
	=
	\int_0^t { \begin{pmatrix}
	1 & 0
	\\0 & 1
	\end{pmatrix}
	\begin{pmatrix}
	d B_s
	\\ \tilde{d B_s}
	\end{pmatrix} }
	\]
	and applied Itos formula to $f(t,x,y) = e^{-t}x+ \sqrt{1-e^{-2t}}y$ which gives us 
	\begin{align*}
	\hat{B}_s = e^{-t}B_s+ \sqrt{1-e^{-2t}}\tilde{B}_s = \int_0^t \frac{e^{-2s}}{\sqrt{1-e^{-2s}}}\tilde{B}_s - e^{-s}B_s \: ds + \int_0^t e^{-s} \: dB_s + \int_0^t \sqrt{1-e^{-2s}} \: d\tilde{B}_s
	\end{align*}
	However I am not sure if this step is necessary. 

	\bigskip
	\textcolor{magenta}{I will consider this problem by Friday.} 
\end{NickQuestion}


\end{document}
